\section{Related Work}

\subsection{Color Palette Generation}

Early work on systematic palette construction focused on perceptual
uniformity in color space selection. Zeileis et al.~\cite{zeileis2009}
proposed HCL-based palettes for statistical graphics, emphasizing
perceptual linearity over aesthetics. Gramazio et al.~\cite{gramazio2017}
introduced \emph{Colorgorical}, a palette generation model that jointly
optimizes perceptual discriminability and aesthetic preference, establishing
discriminability and preference as dual axes of palette quality.
Stone et al.~\cite{stone2014} examined how step count and lightness range
interact in sequential palettes, an observation that motivates our
analysis of even-step advantages in contrast span.

\subsection{Color in Design Systems}

Norman et al.~\cite{norman2024} conducted a qualitative study of
professional designer workflows, identifying three interlinked design spaces:
purpose, palette, and prototype. Their findings motivate computational
scoring that targets purpose-driven constraints (accessibility) rather than
aesthetic preference alone. Color Builder~\cite{colorbuilder2019} provided
a direct-manipulation interface for theme construction but offered no
quantitative quality criteria for the resulting palettes.

\subsection{Color Theme Evaluation}

Yang et al.~\cite{yang2024} proposed a computational model for scoring
color themes through user preference modelling (ACM TAP, 2024).
While the most closely related work in spirit, their approach is
\emph{subjective}—trained on perceptual ratings—and evaluates
multi-color palettes as sets. By contrast, our framework is
\emph{objective}, grounded in geometric properties of the continuous
parametric curve underlying a single-hue ramp, and targets design-system
scalability rather than global aesthetic appeal.

\subsection{Accessibility Metrics}

Existing accessibility evaluation follows WCAG~2.1~\cite{wcag21} contrast
ratio thresholds, which are applied pairwise and do not characterize
palette-level accessibility structure. Prior tools check compliance
post-hoc rather than guiding generation. Our contrast efficiency metric
$\eta$ reformulates accessibility as a structural property of the ramp,
quantifying how economically the required contrast separation is achieved
across the full step range.

\subsection{Positioning}

We are not aware of prior work that addresses the problem of scoring a
\emph{single base color's capacity} to generate an accessible, perceptually
coherent sequential scale across industry-standard step counts---a requirement
specific to design token systems. The proposed framework aims to fill this gap
by combining perceptual geometry, accessibility constraints, and
continuous-curve evaluation into a unified, steps-independent score.
